\documentclass[10pt,a4paper]{amsart}
\usepackage{amsmath,amssymb,amsthm,graphicx}
\usepackage{verbatim}
\newenvironment{reference}{\comment}{\endcomment}
%\newenvironment{reference}{}{}
\newenvironment{slogan}{\comment}{\endcomment}
\newenvironment{history}{\comment}{\endcomment}

% For commutative diagrams we use Xy-pic
\usepackage[all]{xy}

% We use 2cell for 2-commutative diagrams.
\xyoption{2cell}
% Theorem environments.
%
\theoremstyle{plain}
\newtheorem{theorem}[subsection]{Theorem}
\newtheorem{proposition}[subsection]{Proposition}
\newtheorem{lemma}[subsection]{Lemma}

\theoremstyle{definition}
\newtheorem{definition}[subsection]{Definition}
\newtheorem{example}[subsection]{Example}
\newtheorem{exercise}[subsection]{Exercise}
\newtheorem{situation}[subsection]{Situation}

\theoremstyle{remark}
\newtheorem{remark}[subsection]{Remark}
\newtheorem{remarks}[subsection]{Remarks}

\numberwithin{equation}{subsection}

% Macros
%
\def\lim{\mathop{\mathrm{lim}}\nolimits}
\def\colim{\mathop{\mathrm{colim}}\nolimits}
\def\Spec{\mathop{\mathrm{Spec}}}
\def\Hom{\mathop{\mathrm{Hom}}\nolimits}
\def\Ext{\mathop{\mathrm{Ext}}\nolimits}
\def\SheafHom{\mathop{\mathcal{H}\!\mathit{om}}\nolimits}
\def\SheafExt{\mathop{\mathcal{E}\!\mathit{xt}}\nolimits}
\def\Sch{\mathit{Sch}}
\def\Mor{\mathop{\mathrm{Mor}}\nolimits}
\def\Ob{\mathop{\mathrm{Ob}}\nolimits}
\def\Sh{\mathop{\mathit{Sh}}\nolimits}
\def\NL{\mathop{N\!L}\nolimits}
\def\CH{\mathop{\mathrm{CH}}\nolimits}
\def\proetale{{pro\text{-}\acute{e}tale}}
\def\etale{{\acute{e}tale}}
\def\QCoh{\mathit{QCoh}}
\def\Ker{\mathop{\mathrm{Ker}}}
\def\Im{\mathop{\mathrm{Im}}}
\def\Coker{\mathop{\mathrm{Coker}}}
\def\Coim{\mathop{\mathrm{Coim}}}

% Boxtimes
%
\DeclareMathSymbol{\boxtimes}{\mathbin}{AMSa}{"02}

%
% Macros for moduli stacks/spaces
%
\def\QCohstack{\mathcal{QC}\!\mathit{oh}}
\def\Cohstack{\mathcal{C}\!\mathit{oh}}
\def\Spacesstack{\mathcal{S}\!\mathit{paces}}
\def\Quotfunctor{\mathrm{Quot}}
\def\Hilbfunctor{\mathrm{Hilb}}
\def\Curvesstack{\mathcal{C}\!\mathit{urves}}
\def\Polarizedstack{\mathcal{P}\!\mathit{olarized}}
\def\Complexesstack{\mathcal{C}\!\mathit{omplexes}}
% \Pic is the operator that assigns to X its picard group, usage \Pic(X)
% \Picardstack_{X/B} denotes the Picard stack of X over B
% \Picardfunctor_{X/B} denotes the Picard functor of X over B
\def\Pic{\mathop{\mathrm{Pic}}\nolimits}
\def\Picardstack{\mathcal{P}\!\mathit{ic}}
\def\Picardfunctor{\mathrm{Pic}}
\def\Deformationcategory{\mathcal{D}\!\mathit{ef}}
\usepackage{cleveref}
\crefname{originallemma}{original lemma}{original lemmas}
\Crefname{originallemma}{Original lemma}{Original lemmas}
\makeatletter
\newlabel{algebra-lemma-NAK}{{00DV}{source}{}{sourcepointer-1}{}}
\newlabel{algebra-lemma-NAK@cref}{{[originallemma][0][]00DV}{[1][0][]source}{}{}{}}
\makeatother
\crefname{originallemma}{original lemma}{original lemmas}
\Crefname{originallemma}{Original lemma}{Original lemmas}
\makeatletter
\newlabel{decent-spaces-lemma-monomorphism-toward-disjoint-union-dim-0-rings}{{06RZ}{source}{}{sourcepointer-2}{}}
\newlabel{decent-spaces-lemma-monomorphism-toward-disjoint-union-dim-0-rings@cref}{{[originallemma][0][]06RZ}{[1][0][]source}{}{}{}}
\makeatother
\crefname{originalreference}{original reference}{original references}
\Crefname{originalreference}{Original reference}{Original references}
\makeatletter
\newlabel{lemma-locally-quasi-finite-over-field}{{06S0}{source}{}{sourcepointer-3}{}}
\newlabel{lemma-locally-quasi-finite-over-field@cref}{{[originalreference][0][]06S0}{[1][0][]source}{}{}{}}
\makeatother
\title{A monomorphism towards a locally quasi-finite space over a field is locally quasi-finite}
\author{The Stacks Project Authors}
\date{stacks-project:spaces-over-fields.tex:local-snapshot-20261012; GNU FDL 1.2 or later}
\begin{document}
\maketitle
\noindent\textit{Editorial note.} All delivered mathematical text is the author's own, copied from the pinned Stacks Project chapter file spaces-over-fields.tex; no statement or proof step was rewritten, reordered or removed.\par\medskip
\noindent\textit{Editorial note.} The counted proof is the author's complete argument: the decomposition of the target into Artinian pieces and the reduction to the preimages of those pieces, the analysis of a one-point preimage through the epimorphism of zero dimensional local rings, and the application of Nakayama's lemma to obtain finiteness over the field.\par\medskip
\noindent\textit{Editorial note.} The label lemma-mono-towards-locally-quasi-finite-over-field names the counted statement, and the proof environment immediately following it in the source is delivered verbatim as its proof.\par\medskip
\noindent\textit{Editorial note.} The counted statement is carried by the author's own lemma environment in the chapter source.\par\medskip
\noindent\textit{Editorial note.} Pre-open gate: python operations/agent-authoring/tools/precard\_gate.py --source spaces-over-fields.tex --label lemma-mono-towards-locally-quasi-finite-over-field --start 2061 --end 2086 --local-id 28064 returned RESULT=OK (L28064 free, occupancy FREE, public interval clear). Candidate taken from sources/stacks-spaces-over-fields/candidates-halfA.json.\par\medskip
\noindent\textit{Editorial note.} The statement and proof are copied from the pinned chapter file; the cited inputs are kept as external pointers with their canonical Stacks tags.\par\medskip
\noindent\textit{Allowed standard prerequisites (editorial).}\begin{itemize}
\item Structure of a locally quasi-finite space over a field. Same chapter, the lemma describing a locally quasi-finite algebraic space over a field as a disjoint union of spectra of Artinian local algebras finite over the field. Supplies the decomposition of the target into the pieces Spec(A\_i).
\item Monomorphisms towards a disjoint union of zero dimensional rings are schemes. Decent Spaces, the lemma that the source of a monomorphism towards such a disjoint union is a scheme. Supplies the conclusion that X is a scheme, so that the preimages of the pieces can be treated as affine schemes.
\item Nakayama's lemma. Algebra, the lemma allowing the surjectivity of a map of modules to be checked modulo a nilpotent ideal. Supplies the surjectivity of the ring map from the Artinian local ring to the local ring of the preimage from the equality of residue fields.
\end{itemize}\medskip
\ifdefined\hypertarget\hypertarget{sourcepointer-1}{}\fi
\noindent\textit{Source reference.} Original lemma 00DV: sources/stacks-chow/tags.txt, source lines 512--512. Prerequisite pointer invoked at proof line(s) 2081; it names a standard fact of a sibling chapter that the counted argument uses, and every step of the counted proof built on it is retained verbatim in the delivered body.\par\medskip
\ifdefined\hypertarget\hypertarget{sourcepointer-2}{}\fi
\noindent\textit{Source reference.} Original lemma 06RZ: sources/stacks-chow/tags.txt, source lines 8658--8658. Prerequisite pointer invoked at proof line(s) 2073; it names a standard fact of a sibling chapter that the counted argument uses, and every step of the counted proof built on it is retained verbatim in the delivered body.\par\medskip
\ifdefined\hypertarget\hypertarget{sourcepointer-3}{}\fi
\noindent\textit{Source reference.} Original reference 06S0: sources/stacks-spaces-over-fields/pinned-original-source.tex, source lines 2034--2034. Prerequisite pointer invoked at proof line(s) 2069; it names a standard fact of the same chapter that the counted argument uses, and every step of the counted proof built on it is retained verbatim in the delivered body.\par\medskip
\begin{lemma}
\label{lemma-mono-towards-locally-quasi-finite-over-field}
Let $k$ be a field. Let $f : X \to Y$ be a monomorphism of algebraic spaces
over $k$. If $Y$ is locally quasi-finite over $k$ so is $X$.
\end{lemma}

\begin{proof}
Assume $Y$ is locally quasi-finite over $k$. By
Lemma \ref{lemma-locally-quasi-finite-over-field}
we see that $Y = \coprod \Spec(A_i)$ where each $A_i$ is an
Artinian local ring finite over $k$. By
Decent Spaces, Lemma
\ref{decent-spaces-lemma-monomorphism-toward-disjoint-union-dim-0-rings}
we see that $X$ is a scheme. Consider $X_i = f^{-1}(\Spec(A_i))$.
Then $X_i$ has either one or zero points. If $X_i$ has zero points there
is nothing to prove. If $X_i$ has one point, then
$X_i = \Spec(B_i)$ with $B_i$ a zero dimensional local ring
and $A_i \to B_i$ is an epimorphism of rings. In particular
$A_i/\mathfrak m_{A_i} = B_i/\mathfrak m_{A_i}B_i$ and we see that
$A_i \to B_i$ is surjective by Nakayama's lemma,
Algebra, Lemma \ref{algebra-lemma-NAK}
(because $\mathfrak m_{A_i}$ is a nilpotent ideal!).
Thus $B_i$ is a finite local $k$-algebra, and we conclude by
Lemma \ref{lemma-locally-quasi-finite-over-field}
that $X \to \Spec(k)$ is locally quasi-finite.
\end{proof}

\end{document}
